% ----------------------------------------------------------

% Pacotes básicos 
%\usepackage{lmodern}			% Usa a fonte Latin Modern			
\usepackage[T1]{fontenc}		% Selecao de codigos de fonte.
\usepackage[utf8]{inputenc}		% Codificacao do documento (conversão automática dos acentos)
\usepackage{lastpage}			% Usado pela Ficha catalográfica
\usepackage{indentfirst}		% Indenta o primeiro parágrafo de cada seção.
\usepackage{color}		    	% Controle das cores
\usepackage{graphicx}			% Inclusão de gráficos
\usepackage{microtype} 			% para melhorias de justificação
\usepackage{ifthen}		    	% para montar condicionais
\usepackage[brazil]{babel}		% para utilizar termos em portugues
\usepackage[final]{pdfpages}    % para incluir páginas de arquivos pdf
\usepackage{lipsum}				% para geração de dummy text
\usepackage{csquotes}
\usepackage{textcase}           % robust uppercasing
\usepackage{hyperref}
\addto\extrasbrazil{
  \def\equationautorefname{Equação}
  \def\figureautorefname{Figura}
  \def\tableautorefname{Tabela}
  \def\sectionautorefname{Seção}
  \def\subsectionautorefname{Subseção}
}

% Ensure TOC and headers use robust uppercase
\let\MakeUppercase\MakeTextUppercase
\let\MakeLowercase\MakeTextLowercase

%\usepackage[style=long]{glossaries}
%\usepackage{abntex2glossaries}

\usepackage{cancel}         % permite representar o cancelamento de termos em texto ou equacoes

% \usepackage{xcolor} 	    % já incluído em UFPR.sty ou abntex2
\usepackage{smartdiagram}   % gera diagramas a partir de listas
\usepackage{float} 		    % Para a figura ficar na posição correta
\usepackage{textcomp} 		% supporte para fontes da Text Companion 
\usepackage{longtable}		% uso de longtable
\usepackage{amsmath}	    % simbolos matematicos
\usepackage{lscape}         % páginas em paisagem
\usepackage{multicol}       % mescla de colunas em tabelas
\usepackage{multirow}       % mescla de linhas em tabelas
\usepackage{newfloat}       % criação do indice de quadros
\usepackage{caption} 		% configura legenda
%[format=plain]
	%\renewcommand\caption[1]{%
    \captionsetup{font=small}	% tamanho da fonte 10pt
    %,format=hang
 	% \caption{#1}}
	%\captionsetup{width=0.8\textwidth}

\usepackage{booktabs}
\usepackage{adjustbox}
\usepackage{rotating}
\usepackage{array}
\usepackage{listings}
\usepackage[linesnumbered,ruled,vlined,portuguese]{algorithm2e}

% \usepackage{xcolor}

\SetKwInput{Entrada}{Entrada}
\SetKwInput{Saida}{Saída}
\SetKwFor{Para}{para}{faça}{fim para}
\SetKwFor{Enquanto}{enquanto}{faça}{fim enquanto}
\SetKw{Retorne}{retorne}
\SetKwIF{If}{ElseIf}{Else}{se}{então}{senão se}{senão}{fim se}

% Designação e lista própria dos algoritmos, conforme a NBR 14724
\SetAlgorithmName{ALGORITMO}{algoritmo}{LISTA DE ALGORITMOS}
\SetAlgoCaptionSeparator{~\textendash~}

% Designação e lista própria das listagens de código, conforme a NBR 14724
\renewcommand{\lstlistingname}{PROGRAMA}
\renewcommand{\lstlistlistingname}{LISTA DE PROGRAMAS}

% Numeração contínua, à semelhança das figuras, das tabelas e dos algoritmos.
% O contador é criado pelo pacote listings apenas no início do documento.
\AtBeginDocument{%
    \counterwithout{lstlisting}{chapter}%
    \renewcommand{\thelstlisting}{\arabic{lstlisting}}%
}

% Nas listas próprias, cada item é designado por seu nome específico, como já
% ocorre com as figuras e com as tabelas
\makeatletter
\renewcommand*{\l@algocf}[2]{%
    \begingroup
        \renewcommand{\numberline}[1]{%
            \hb@xt@\@tempdima{\algorithmcfname~##1\hfill\textendash\hfill}}%
        \@dottedtocline{1}{0pt}{3.6cm}{#1}{#2}%
    \endgroup}
\renewcommand*{\l@lstlisting}[2]{%
    \begingroup
        \renewcommand{\numberline}[1]{%
            \hb@xt@\@tempdima{\lstlistingname~##1\hfill\textendash\hfill}}%
        \@dottedtocline{1}{0pt}{3.6cm}{#1}{#2}%
    \endgroup}
\makeatother

\lstset{
    basicstyle=\ttfamily\small,
    keywordstyle=\color{blue},
    stringstyle=\color{red},
    commentstyle=\color{green!60!black},
    breaklines=true,
    frame=single,
    numbers=left,
    numberstyle=\tiny,
    captionpos=t,
    showstringspaces=false,
    keepspaces=true,
    breakatwhitespace=false,
    aboveskip=0.8em,
    belowskip=0.8em
}

% ----------------------------------------------------------

% Indicação da fonte das ilustrações e das tabelas, obrigatória pela NBR 14724,
% mesmo quando o material é de elaboração do próprio autor.
% O comando homônimo do abnTeX2 é redefinido para adotar a grafia em caixa alta
% e o alinhamento à esquerda exigidos pelas normas da UFPR.
\renewcommand{\fonte}[1]{%
    \par\vspace{1mm}%
    \noindent{\footnotesize\raggedright FONTE:~#1\par}%
}

% Nota explicativa opcional, posicionada abaixo da fonte.
\renewcommand{\nota}[1]{%
    \par\noindent{\footnotesize\raggedright NOTA:~#1\par}%
}

% ----------------------------------------------------------

% Pacotes de citações BibLaTeX
% A NBR 10520 exige que citações simultâneas de vários documentos sejam
% apresentadas em ordem alfabética, o que é obtido com a opção sortcites.
% A NBR 6023 não prevê informações acrescentadas às referências, razão pela
% qual a remissão às páginas de citação (backref) permanece desativada.
\usepackage[style=abnt,backref=false,backend=biber,sortcites=true,language=brazilian]{biblatex}

% Forçar nomes de autores em maiúsculas e minúsculas (Capitalizado)
\renewcommand{\mkbibnamefamily}[1]{#1}
\renewcommand{\mkbibnamegiven}[1]{#1}
\renewcommand{\mkbibnameprefix}[1]{#1}
\renewcommand{\mkbibnamesuffix}[1]{#1}

\AtBeginBibliography{
  \renewcommand{\mkbibnamefamily}[1]{\MakeUppercase{#1}}
  \renewcommand{\mkbibnamegiven}[1]{#1}
  \renewcommand{\mkbibnameprefix}[1]{#1}
  \renewcommand{\mkbibnamesuffix}[1]{#1}
}


% ----------------------------------------------------------

% alterando o aspecto da cor azul
\definecolor{blue}{RGB}{55,10,249}

% ----------------------------------------------------------

\PrepareListOf{quadro}{
    \renewcommand{\cftfigpresnum}{Quadro~}
}

\DeclareFloatingEnvironment[
    fileext=loq,
    listname={\textbf{LISTA DE QUADROS}},
    name=Quadro,
    %placement=p,
    within= none, % numeracao continua
    %within=section, % numeracao reinicia em cada seccao
    %chapterlistsgaps=off
]{quadro}

\newlistentry{quadro}{loq}{0}

% Customize ‘List of Diagrams’
\PrepareListOf{quadro}{
    \renewcommand{\cftquadropresnum}{\normalsize{QUADRO}~}
    \setlength{\cftquadronumwidth}{3.2cm}
    %\renewcommand{\cftquadroname}{\quadroname\space} 
    \renewcommand*{\cftquadroaftersnum}{\hfill--\hfill}
}

\makeatletter
%% we define a helper macro for adjusting lists of new floats to
%% accept a * behind them for not being shown in the TOC, like
%% the other list printing commands in memoir
\newcommand{\AdjustForMemoir}[1]{
  \csletcs{kept@listof#1}{listof#1}
  \csdef{listof#1}{%
    \@ifstar
     {\csappto{newfloat@listof#1@hook}{\append@star}%
      \csuse{kept@listof#1}}%
     {\csuse{kept@listof#1}}%
  }
}

\def\append@star#1{#1*}
\makeatother

\AdjustForMemoir{quadro} % prepare `\listofdirfigures` so it accepts a *

\makeatletter
\let\oldcontentsline\contentsline
\def\contentsline#1#2{%
    \expandafter\ifx\csname l@#1\endcsname\l@section
	\expandafter\@firstoftwo
	\else
	\expandafter\@secondoftwo
	\fi
	{%
		\oldcontentsline{#1}{#2}%
	}{%
	\normalsize %ajusta tamanho da fonte na lista
	\oldcontentsline{#1}{#2}%
}%
}
\makeatother

% ----------------------------------------------------------

% Ajusta indentação de Referencias no ToC
\defbibheading{bay}[\bibname]{
  \chapter*{#1}%
  \markboth{#1}{#1}%
  \addcontentsline{toc}{chapter}
  {\protect\numberline{}\bibname}
}

\makeatletter
\pretocmd{\chapter}{\addtocontents{toc}{\protect\addvspace{-5\p@}}}{}{}
\pretocmd{\section}{\addtocontents{toc}{\protect\addvspace{2\p@}}}{}{}
\makeatother

% ----------------------------------------------------------

% Conforme a NBR 6027, o sumário deve reproduzir os títulos tal como aparecem no
% texto. Como as seções primárias e secundárias são impressas em caixa alta, as
% entradas correspondentes do sumário recebem o mesmo tratamento.
\makeatletter
\settocpreprocessor{chapter}{%
    \let\ABNTEXchaptertoc\f@rtoc%
    \def\f@rtoc{%
        \texorpdfstring{\bfseries\MakeTextUppercase{\ABNTEXchaptertoc}}{\ABNTEXchaptertoc}}%
}
\makeatother

% O memoir só oferece pré-processador de sumário para livro, parte e capítulo.
% Para as seções secundárias, a conversão é aplicada no momento em que a entrada
% é gravada no arquivo de sumário.
\let\ABNTEXaddcontentsline\addcontentsline
\renewcommand{\addcontentsline}[3]{%
    \ifthenelse{\equal{#2}{section}}%
        {\ABNTEXaddcontentsline{#1}{#2}{\protect\MakeTextUppercase{#3}}}%
        {\ABNTEXaddcontentsline{#1}{#2}{#3}}%
}
